%% =====================================================================
%%  Adani Ports and Special Economic Zone Ltd (APSEZ): A Strategic Analysis
%%  Standalone LaTeX source. Compile with pdfLaTeX (run twice so that the
%%  table of contents and cross-references resolve).
%% =====================================================================
\documentclass[12pt, a4paper]{article}
\usepackage[utf8]{inputenc}
\usepackage[T1]{fontenc}
\usepackage{lmodern}
\usepackage[british]{babel}
\usepackage{geometry}
\geometry{a4paper, margin=1in}
\usepackage{setspace}
\usepackage{titlesec}
\usepackage{enumitem}
\usepackage{booktabs} % professional table rules
\usepackage{caption}
\usepackage{graphicx}
\usepackage{float}
\usepackage[section]{placeins}

% Additional packages needed by this report's content
\usepackage{amsmath}
\usepackage{textcomp}
\usepackage{microtype}
\usepackage{tfrupee}        % Indian rupee sign (\rupee)
\usepackage{pifont}         % check and cross marks
\usepackage{newunicodechar}
\usepackage{array}
\usepackage{tabularx}
\usepackage{xltabular}      % full-width tables that may break across pages
\usepackage{needspace}
\usepackage{xcolor}
\usepackage{tikz}           % diagrams (Exhibits 3 and 10)
\usetikzlibrary{positioning,arrows.meta}
\usepackage{xurl}

\usepackage{hyperref}
\hypersetup{
colorlinks=true,
linkcolor=blue,
filecolor=magenta,
urlcolor=cyan,
pdftitle={Adani Ports and Special Economic Zone Ltd (APSEZ): A Strategic Analysis},
}

%% ---------------------------------------------------------------------
%%  Exhibits: numbered floats ("Exhibit 1", "Exhibit 2", ...) with the
%%  caption above the content, styled like standard table captions
%% ---------------------------------------------------------------------
\floatstyle{plaintop}
\newfloat{exhibit}{htbp}{loe}
\floatname{exhibit}{Exhibit}
\captionsetup[exhibit]{position=top}

%% ---------------------------------------------------------------------
%%  Diagram colours (Exhibits 3 and 10 refer to colour-coded ratings)
%% ---------------------------------------------------------------------
\definecolor{muted}{HTML}{5B6470}
\definecolor{rulegrey}{HTML}{C9CED6}
\definecolor{soft}{HTML}{F2F4F7}
\definecolor{accent}{HTML}{1F4E79}
\definecolor{lowbg}{HTML}{DCEFE3}
\definecolor{lowink}{HTML}{1E5E3A}
\definecolor{modbg}{HTML}{FBECCC}
\definecolor{modink}{HTML}{7A5300}
\definecolor{highbg}{HTML}{F6D5D2}
\definecolor{highink}{HTML}{8A2219}

%% ---------------------------------------------------------------------
%%  Helper macros
%% ---------------------------------------------------------------------
% Source note beneath an exhibit
\newcommand{\exsource}[1]{\par\vspace{4pt}{\footnotesize\raggedright #1\par}}
% Assessment symbols
\DeclareRobustCommand{\halfcirc}{\tikz[baseline=-0.62ex]{%
  \draw[line width=0.45pt] (0,0) circle (0.78ex);
  \fill (0,0.78ex) arc (90:270:0.78ex) -- cycle;}}
\DeclareRobustCommand{\yes}{\ding{51}}
\DeclareRobustCommand{\no}{\ding{55}}
\DeclareRobustCommand{\partly}{\halfcirc}
% Proposals shown inside the Ansoff matrix
\newcommand{\prop}[1]{\textit{#1}}

% Unicode characters used in the text
\newunicodechar{₹}{{\rupee}}
\DeclareUnicodeCharacter{2192}{\ensuremath{\rightarrow}}% →
\newunicodechar{≈}{\ensuremath{\approx}}
\newunicodechar{≥}{\ensuremath{\geq}}
\newunicodechar{−}{\ensuremath{-}}
\newunicodechar{⅓}{\ensuremath{\tfrac{1}{3}}}
\newunicodechar{✓}{\ding{51}}
\newunicodechar{✗}{\ding{55}}
\newunicodechar{◐}{\halfcirc}

% Column types
\newcolumntype{L}{>{\raggedright\arraybackslash}X}
\newcolumntype{R}{>{\raggedleft\arraybackslash}X}
\newcolumntype{P}[1]{>{\raggedright\arraybackslash}p{#1}}
\newcolumntype{Q}[1]{>{\raggedleft\arraybackslash}p{#1}}
\newcolumntype{S}{>{\centering\arraybackslash}p{5.4mm}}

% Common table set-up (booktabs rules, compact text so tables fit the margins)
\newcommand{\tablesetup}{\footnotesize\renewcommand{\arraystretch}{1.25}\setlength{\tabcolsep}{4pt}}
\newcommand{\rowsep}{\addlinespace[3pt]}
\newcommand{\reclabel}[1]{\textbf{#1}}

% Compact lists inside table cells
\newlist{celllist}{itemize}{1}
\setlist[celllist]{label=\textbullet, leftmargin=1.0em, itemsep=0.2ex,
  topsep=0.2ex, parsep=0pt, partopsep=0pt, before=\vspace{-0.3ex}, after=\vspace{-0.9ex}}

%% ---------------------------------------------------------------------
%%  Title
%% ---------------------------------------------------------------------
\title{
\vspace*{2cm}
\Huge \textbf{Adani Ports and Special Economic Zone Ltd (APSEZ): A Strategic Analysis} \\
\vspace{0.5cm}
\Large Competitive position, resources and growth strategy assessed with Porter's Five Forces, VRIN, VRIO and the Ansoff matrix
}

\author{
\textbf{Assignment:} SCM Assignment: Strategic analysis of APSEZ \\
\textbf{Student:} [Name / ID] \\
}
\date{\textbf{Date:} November 2026}

\begin{document}
%% =====================================================================

%% ---------------------------------------------------------------------
%%  Title page
%% ---------------------------------------------------------------------
\maketitle
\thispagestyle{empty}
\vspace{2cm}
\begin{center}
\renewcommand{\arraystretch}{1.4}
\begin{tabular}{@{}lp{10cm}@{}}
\toprule
Evidence cut-off & 28 October 2026 (Q2 FY27 results) \\
Referencing      & Harvard (Cite Them Right) \\
Word count       & 4,977 words (main body; excludes exhibits, references and appendices) \\
\bottomrule
\end{tabular}
\end{center}
\newpage

\pagenumbering{roman}
\tableofcontents
\vspace{1cm}
\noindent{\footnotesize
Units: ₹ cr = ₹ crore (₹10 million). MMT = million metric tonnes. MTPA =
million tonnes per annum. TEU = twenty-foot equivalent unit. FY = Indian
financial year ending 31 March (FY26 = April 2025--March 2026). Q2 FY27 =
July--September 2026. ROCE = return on capital employed, as defined by APSEZ
unless stated otherwise.\par}
\newpage
\pagenumbering{arabic}

%% =====================================================================
\section{Executive summary}\label{sec:summary}
%% =====================================================================
APSEZ is India's largest commercial port operator. It handled 500.8 MMT of
cargo in FY26 and reports a 27.1\% share of national port volume (APSEZ,
2026a). This report tests that position against four frameworks. The evidence
base is the FY26 audited accounts and disclosures up to 28 October 2026.

The analysis reaches four conclusions. First, Indian port services are a
moderately attractive industry for incumbents with good sites. In containers,
however, rivalry and buyer power are high and rising, and a state-led port at
Vadhavan is planned to add large west-coast capacity from around 2030, possibly
later. Second, APSEZ has one clear sustained advantage: the Mundra complex. That
complex combines a port concession, contiguous SEZ land, deep water and terminal
joint ventures with two of the world's largest container lines, and it carries
45\% of APSEZ's domestic volume (APSEZ, 2026d). Most other resources, including
scale, the logistics network and access to capital, deliver parity or a
temporary advantage. Third, the Mundra advantage is legally contingent. The
concession, and the container sub-concessions tied to it, expire on 16 February
2031, and several supplementary agreements remain unsigned (APSEZ, 2026b).
Fourth, new capital is flowing into segments where APSEZ has parity and earns
8--13\% ROCE, against 23\% in domestic ports. Diversification therefore spreads
risk but dilutes returns.

Four recommendations follow:
\begin{itemize}
  \item secure Mundra's tenure before 2029, and make long-payback investment at
        Mundra conditional on that tenure;
  \item introduce a return gate for logistics, marine and international investment;
  \item defend container share on the west coast before Vadhavan opens;
  \item turn around or recycle the under-performing non-Mundra ports.
\end{itemize}

%% =====================================================================
\section{Introduction and method}\label{sec:method}
%% =====================================================================
The report evaluates APSEZ's competitive position, resources and growth
strategy. The analysis runs from the outside in. Porter's Five Forces is
applied first, at industry level, so that APSEZ's own strengths do not shape
the industry judgement. VRIN (Barney, 1991) and VRIO (Barney, 1995; Barney and
Hesterly, 2021) then assess the firm's resources. The Ansoff (1957) matrix
classifies growth initiatives, which are evaluated for suitability,
acceptability and feasibility (Whittington \textit{et al.}, 2023). The final
step reconciles the four sets of findings and derives recommendations.

Evidence comes from exchange filings, audited financial statements and
government data wherever possible. Company claims are identified as such. Press
articles that reproduce company releases are not treated as independent
corroboration. All derived ratios are the author's calculations, and their
scope is stated, because consolidated, segment and joint-venture figures are not
interchangeable. \hyperref[app:data]{Appendix~B} lists the data limitations.

%% =====================================================================
\section{Company overview and financial profile}\label{sec:company}
%% =====================================================================
APSEZ was incorporated in 1998. Its foundation is a 30-year concession for
Mundra port in Gujarat, granted by the Government of Gujarat and the Gujarat
Maritime Board (GMB) and effective from 17 February 2001 (APSEZ, 2026b). Growth
came first from building out Mundra and its special economic zone, and later
from acquisitions and new concessions on every Indian coastline. APSEZ now
operates 15 domestic ports and terminals, plus four international ports: Haifa
(Israel), Colombo West International Terminal (Sri Lanka), Dar es Salaam
(Tanzania) and the North Queensland Export Terminal (NQXT, Australia). Domestic
capacity reached 671 MTPA in September 2026 after APSEZ won a 30-year concession
for two dry-bulk berths at Paradip (APSEZ, 2026d).

The business has four reporting segments: domestic ports, international ports,
logistics, and marine services. A small SEZ segment sits alongside them.
Exhibits~\ref{ex:profile} and~\ref{ex:segments} summarise the financial profile,
which shows four features.
\begin{itemize}
  \item \textbf{Growth has been fast, but margins have been flat.} Revenue grew
        at 22.6\% a year between FY22 and FY26, while the EBITDA margin stayed
        between 59\% and 62\%. Over that period the business mix moved towards
        lower-margin activities.
  \item \textbf{Earnings are concentrated.} Domestic ports produced two-thirds
        of FY26 revenue and EBITDA equal to about 82\% of the consolidated total,
        at a 73\% margin and 23\% ROCE. Within that segment, Mundra alone handled
        45\% of domestic volume in H1 FY27. Volume at the other domestic ports
        was flat year on year (APSEZ, 2026d).
  \item \textbf{The business is capital-intensive.} Gross fixed assets are about
        4.1 times revenue. FY26 capex of ₹15,320 cr was about 40\% of revenue and
        exceeded guidance by 28--39\% (APSEZ, 2026a). Leverage is moderate: net
        debt was 1.9 times EBITDA at the end of FY26.
  \item \textbf{Recent growth is partly bought.} International volume rose from
        17 to 47 MMT in H1 FY27, mostly because of NQXT. APSEZ bought NQXT in
        December 2025 from a promoter-group company and paid in new shares
        (APSEZ, 2025). Domestic volume grew only 2.6\% in the same period.
\end{itemize}

Two further facts matter for later sections. In FY26, 16.1\% of revenue came
from Adani group companies, up from 11.9\% in FY25. JVs added another 3.3\%
(APSEZ, 2026b, note 32; consistent with the 19.8\% related-party share reported
in APSEZ, 2026c). Separately, no single external customer accounts for 10\% or
more of revenue.

% ---------------------------------------------------------------- Exhibit 1
\begin{exhibit}[!htbp]
\caption{APSEZ consolidated profile, FY22--FY26 and H1 FY27}\label{ex:profile}
{\tablesetup
\begin{tabularx}{\linewidth}{P{41mm}RRRRR}
\toprule
\textbf{₹ cr unless stated} & \textbf{FY22} & \textbf{FY24} & \textbf{FY25} & \textbf{FY26} & \textbf{H1 FY27} \\
\midrule
Revenue from operations & 17,119 & 26,711 & 31,079 & 38,736 & 22,161 \\ \rowsep
EBITDA (margin) & 10,397 (61\%) & 15,864 (59\%) & 19,025 (61\%) & 22,851 (59\%) & 13,347 (60\%) \\ \rowsep
Profit after tax & 4,953 & 8,104 & 11,061 & 12,782 & 7,332 \\ \rowsep
Gross fixed assets (excl.\ CWIP) & 74,076 & 93,006 & 1,11,163 & 1,58,850 & n/d \\ \rowsep
Capex (vs guidance) & n/d & n/d & n/d & 15,320 (11--12k) & 10,521 (FY27 G: 15--17k) \\ \rowsep
Net debt / EBITDA (x) & n/d & n/d & n/d & 1.9 & 1.8 \\ \rowsep
Cargo, MMT (domestic) & --- & --- & 450.2 & 500.8 (≈451) & 280 (233) \\ \rowsep
All-India cargo / container share & --- & --- & 27.0\% / 45.5\% & 27.1\% / 45.5\% & 27.4\% / 44.6\% \\ \rowsep
Consolidated ROCE (company-defined) & --- & 13\% & 15\% & 16\%* & --- \\
\bottomrule
\end{tabularx}}
\exsource{Sources: APSEZ (2026a; 2026d). FY figures are the management
presentation basis. The audited FY25 revenue is ₹30,475 cr (APSEZ, 2026b),
against the restated ₹31,079 cr shown here. *Excludes NQXT. n/d = not disclosed
in the documents reviewed. H1 FY27 market shares are compared with H1 FY26
(28.0\% / 45.5\%).}
\end{exhibit}

% ---------------------------------------------------------------- Exhibit 2
\begin{exhibit}[!htbp]
\caption{Segment economics, FY26 (and Q2 FY27 margin)}\label{ex:segments}
{\tablesetup
\begin{tabularx}{\linewidth}{P{21mm}Q{20mm}Q{16mm}Q{21mm}Q{17mm}L}
\toprule
\textbf{Segment} & \textbf{Revenue ₹ cr (YoY)} & \textbf{EBITDA ₹ cr} & \textbf{EBITDA margin FY26 (Q2 FY27)} & \textbf{ROCE FY26 (FY25)} & \textbf{What drives it} \\
\midrule
Domestic ports & 25,755 (+13\%) & 18,849 & 73\% (72\%) & 23\% (21\%) & Mundra; captive harbour services (90\% margin) \\ \rowsep
International ports & 4,539 (+34\%) & ≈1,298 & 29\% (48\%) & 8\%* (6\%) & NQXT (from Q4 FY26), Colombo ramp-up \\ \rowsep
Logistics & 4,478 (+55\%) & 863 & 19\% (19\%) & 10\% (6\%) & Trucking 8.4\% and freight-forwarding 10.4\% margins (Q2 FY27) \\ \rowsep
Marine & 2,681 (+134\%) & 1,357 & 51\% (44\%) & 13\% (13\%) & Acquired offshore-vessel fleets (Middle East, Africa, South Asia) \\
\bottomrule
\end{tabularx}}
\exsource{Sources: APSEZ (2026a; 2026d). International EBITDA is derived from
the reported 28.6\% margin. *Excludes NQXT. Segment figures are before
consolidation adjustments and SEZ income.}
\end{exhibit}

%% =====================================================================
\section{The Indian ports industry}\label{sec:industry}
%% =====================================================================
Indian ports operate under two regimes:
\begin{itemize}
  \item \textbf{Major ports.} Twelve operating major ports follow a landlord
        model under the Major Port Authorities Act 2021. Private operators such
        as DP World, PSA and APM Terminals run terminals there under
        public--private partnerships. Major ports handled 915.17 MT in FY26, up
        7.1\% (Paradip Port Authority, 2026).
  \item \textbf{Non-major ports.} These are governed by state maritime boards.
        Under the Indian Ports Act 2025, those boards or their concessionaires
        set tariffs (Government of India, 2025). Using APSEZ's estimate of 1,668
        MMT for total Indian traffic, non-major ports handled about 45\% of the
        national total. Gujarat accounts for much of that, and APSEZ is its
        largest operator. In September 2026 Mundra became the only non-major
        port among the first four ports designated ``Mega Ports'' (Maritime
        Gateway, 2026).
\end{itemize}

The regulatory split has economic consequences. Private non-major ports in
Gujarat have long set their own commercial rates, subject to a royalty paid to
the state. Major-port terminals are run under concession agreements whose
tariffs are set by the port authority's board. The 2025 Act moves both systems
towards market-based pricing, with an advisory body overseeing tariff
transparency (Government of India, 2025). For APSEZ this cuts both ways. It
keeps APSEZ's pricing freedom, and it also lets public ports and their terminal
operators compete more flexibly on price.

Demand is moving in different directions by cargo type. Container volumes are
growing. Coal, about 30\% of APSEZ's domestic cargo in Q2 FY27, fell 6\% year on
year (APSEZ, 2026d). Liquids are tied to refineries. Hinterland reach is
widening as the Western Dedicated Freight Corridor carries double-stack
container trains. This is consistent with Notteboom and Rodrigue's (2005)
account of port regionalisation, in which inland networks increasingly decide
which port wins. Empirical work on India links competition between major and
non-major ports to lower efficiency and possible over-supply of infrastructure
(Adler \textit{et al.}, 2022). That is a familiar risk in a sector where
capacity arrives in large, lumpy increments (Haralambides, 2002).

%% =====================================================================
\section{Porter's Five Forces}\label{sec:fiveforces}
%% =====================================================================
The industry is defined as commercial cargo-handling services at Indian gateway
ports. The forces differ sharply by cargo type, so each is rated separately for
containers, dry bulk and liquids (Exhibit~\ref{ex:fiveforces}). Porter (2008)
treats government as an influence on the forces rather than a force in its own
right. One adaptation is made here: the concession grantor is treated as a
supplier, because it supplies the indispensable input, the right to operate a
waterfront, and it sets royalties and renewal terms.

\subsection{Competitive rivalry: high in containers, moderate in bulk}
Container competition on the west coast is intensifying. APSEZ's all-India
container share fell from 45.9\% to 44.4\% between Q2 FY26 and Q2 FY27 (APSEZ,
2026d). Over the same period JNPA, the main public alternative near Mumbai,
passed 8 million TEU in FY26 for the first time (Free Press Journal, 2026).
APSEZ's domestic capacity utilisation is about 69\% (451 MMT handled against 653
MTPA), and JSW Infrastructure plans to reach 400 MTPA by FY30 (Whalesbook,
2026). Spare capacity of this kind encourages price competition (Haralambides,
2002). In bulk, rivalry is regional rather than national. Even so, APSEZ's
east-coast ports saw EBITDA margin fall from 50\% to 41\% year on year in Q2
FY27, which points to price or mix pressure near Paradip and Visakhapatnam.

\subsection{Threat of new entry: moderate}
Building a new port requires a deep-water site, coastal and environmental
clearances, land, and very large capital. APSEZ's gross fixed assets alone are
₹1.59 lakh crore. These barriers are real. However, the most credible entrant
is the state itself. Vadhavan, a ₹76,220 cr port planned by JNPA and
Maharashtra, targets about 23.2 million TEU at full build, with its first phase
due around 2030, although some estimates are later (ThePrint, 2024). Landlord
concessions at major ports also let well-funded operators enter without owning
a port, so the barriers protect existing sites more than they protect the
market.

\subsection{Supplier power: moderate, rising to high for Mundra}
Two suppliers matter most. The first is the concession grantor. GMB prescribes
waterfront royalty rates (APSEZ, 2026b) and will decide the terms on which
Mundra continues after February 2031. APSEZ is reported to have applied for an
extension (India Shipping News, 2021; CargoConnect, n.d.), but no final decision
has been published. The second is Indian Railways, the monopoly provider of
rail haulage, whose tariffs and rake supply shape port catchments. APSEZ's rail
container volume fell 6.4\% year on year in Q2 FY27. Equipment makers and
dredging contractors matter less, partly because APSEZ owns 64 dredgers.

\subsection{Buyer power: high in containers, low to moderate in bulk and liquids}
Global container lines are concentrated and choose where to call. At Mundra,
the two largest container terminals are 50:50 joint ventures with MSC and CMA
CGM (APSEZ, 2026b). This makes the largest buyers partners as well. Such
co-opetition (Brandenburger and Nalebuff, 1996) secures volume today, but it
gives those partners leverage as their sub-concessions approach their 2031
end-date. In bulk and liquids, buyer power is weaker. Plants located near
ports, conveyor connections and refinery jetties create high switching costs,
and 16\% of revenue comes from group companies. No external customer exceeds
10\% of revenue.

A cargo owner's ability to switch also depends on distance. Containers bound
for northern India can be routed through Mundra, Pipavav or JNPA, and the
freight corridor narrows the cost gap between them. A refinery or power plant
tied to a dedicated jetty cannot switch at all. In Porter's terms, buyer power
is therefore highest exactly where APSEZ's growth in market share has stalled.

\subsection{Threat of substitutes: low overall, moderate for coal}
Sea transport has no practical substitute for India's export and import trade.
The real substitution risk lies in the cargo itself. As India's demand for
imported thermal coal declines, coal-handling capacity loses value. Coal is
still the largest single cargo at several APSEZ ports.

\textbf{Industry verdict.} The industry is moderately attractive for incumbents
with good sites. Bulk and liquids are protected by captive demand, and the
container gateway segment is becoming less attractive towards 2030. APSEZ's
72--73\% domestic margin is higher than this structure would predict. Three
facts explain much of the gap. Part of the cargo is captive. Harbour services
such as towage and pilotage, which only the port owner can supply inside its
own port, earn about 90\% margins. And the most contested activity, container
handling at Mundra, sits mostly in joint ventures that are not consolidated.
AICTPL, the MSC joint venture, earned about 57\% in Q2 FY27 (APSEZ, 2026d). Firm
advantage, examined next, therefore has to explain less of the margin than the
headline suggests.

% ---------------------------------------------------------------- Exhibit 3
\begin{exhibit}[H]
\caption{Five Forces for Indian gateway port services, with ratings by segment
(FY26--H1 FY27; ``box'' = container)}\label{ex:fiveforces}
\centering
\begin{tikzpicture}[
    font=\fontsize{7}{8.8}\selectfont,
    execute at begin node={\hyphenpenalty=10000\exhyphenpenalty=10000},
    force/.style={draw=rulegrey, rounded corners=3pt, align=center,
                  inner sep=4pt, line width=0.6pt},
    flow/.style={-{Stealth[length=5pt,width=4pt]}, draw=muted, line width=1pt}]
  \node[force, fill=highbg, text width=53mm] (riv) {%
    \textcolor{highink}{\textbf{RIVALRY: MODERATE–HIGH, RISING}}\\
    Box HIGH · Bulk MOD · Liquid LOW–MOD\\
    APSEZ box share 45.9\% → 44.4\% (Q2)\\
    JNPA >8m TEU (FY26); JSW → 400 MTPA\\
    APSEZ domestic utilisation ≈69\%\\
    East-coast margin 50\% → 41\%};
  \node[force, fill=modbg, text width=62mm, above=9mm of riv] (ent) {%
    \textcolor{modink}{\textbf{THREAT OF NEW ENTRY: MODERATE}}\\
    High greenfield barriers (capital, draft, clearances)\\
    but the state is an entrant: Vadhavan, ₹76,220 cr, c.2030\\
    Landlord PPPs allow entry without owning a port};
  \node[force, fill=modbg, text width=41mm, left=6mm of riv] (sup) {%
    \textcolor{modink}{\textbf{SUPPLIERS: MODERATE}}\\
    \textcolor{highink}{\textbf{→ HIGH for Mundra near 2031}}\\
    Grantor (GMB): royalty, renewal\\
    Indian Railways: haulage monopoly\\
    64 own dredgers limit contractors};
  \node[force, fill=highbg, text width=41mm, right=6mm of riv] (buy) {%
    \textcolor{highink}{\textbf{BUYERS: MODERATE–HIGH}}\\
    Box HIGH · Bulk/Liquid LOW–MOD\\
    MSC, CMA CGM: buyers and JV partners\\
    No external customer ≥10\% of revenue\\
    16\% of revenue from group companies};
  \node[force, fill=lowbg, text width=62mm, below=9mm of riv] (subs) {%
    \textcolor{lowink}{\textbf{SUBSTITUTES: LOW (MODERATE for coal)}}\\
    No practical alternative to sea EXIM trade\\
    Falling thermal-coal imports (APSEZ coal −6\% YoY)};
  \draw[flow] (ent.south) -- (riv.north);
  \draw[flow] (sup.east) -- (riv.west);
  \draw[flow] (buy.west) -- (riv.east);
  \draw[flow] (subs.north) -- (riv.south);
\end{tikzpicture}
\exsource{Sources: APSEZ (2026a; 2026b; 2026d); Free Press Journal (2026);
ThePrint (2024); Whalesbook (2026). Ratings are the author's assessment.}
\end{exhibit}

%% =====================================================================
\section{VRIN analysis}\label{sec:vrin}
%% =====================================================================
Barney (1991) argues that a resource supports a sustained advantage when it is
valuable, rare, imperfectly imitable and non-substitutable.
Non-substitutability is a distinct test. It asks whether a rival could carry
out the same strategy using a \emph{different} bundle of resources, which is not
the same as Porter's question of whether customers can switch to a different
product. Ten resources are assessed in Exhibit~\ref{ex:vrin}.

Each resource is tested on evidence that APSEZ actually holds it, and the
assessment does not assume that something important must be an advantage. Two
resources commonly cited for APSEZ, scale and its integrated network, are
deliberately split into their parts, because the parts behave differently.

\textbf{The Mundra complex (R1)} is the only resource that clearly meets the
first three tests. It is valuable: Mundra earns a 72\% EBITDA margin and carries
45\% of domestic volume. It is rare: by APSEZ's account it is the country's
largest commercial port, it is the only private Mega Port, and it sits next to
a notified SEZ. It is hard to imitate because it was built over 25 years.
Dierickx and Cool (1989) describe this kind of time-compression diseconomy:
land, permits, deep-water berths and an industrial cluster around them cannot
be bought quickly. Non-substitutability is only partial. For containers, a
rival bundle of JNPA terminals, Vadhavan and freight-corridor rail could serve
the same northern hinterland. For bulk and SEZ-linked cargo, no comparable
bundle exists. Imitability also depends on law, not only on economics: much of
R1 is a right that expires in 2031.

\textbf{The non-Mundra portfolio (R2)} fails the imitability test. Most of
these ports were acquired. Barney (1986) shows that resources bought in
competitive factor markets rarely earn above-normal returns unless the buyer
has superior information or integration skill. Here the evidence of superior
integration is weak. Volumes were flat in H1 FY27, east-coast margins fell, and
the Visakhapatnam coal terminal ended in arbitration and is held for sale
(APSEZ, 2026b).

\textbf{Transshipment positions (R3)}, at Vizhinjam (about 1.3 million TEU in
its first full year; Shipping Tribune, 2026) and Colombo, are valuable and
currently rare in India. They are only partly protected, because the state
plans further hubs and Colombo's other terminals and Singapore are close
substitutes. \textbf{Joint ventures with container lines (R4)} are valuable but
less rare than they look, since APM Terminals and DP World hold similar
positions at JNPA and Pipavav. \textbf{Rail and logistics parks (R5)} and
\textbf{trucking and freight forwarding (R6)} fail on rarity: CONCOR and many
private firms offer similar services, and the freight corridor is public
infrastructure. \textbf{The marine fleet (R7)} creates value, but much of it
comes from captive harbour services, a monopoly that derives from owning R1.
\textbf{Financial capacity (R9)} fails non-substitutability, because sovereign
capital funds rival capacity. \textbf{Digital systems (R10)} are bought from
vendors and serve as a parity control.

% ---------------------------------------------------------------- Exhibit 4
\begin{exhibit}[!htbp]
\caption{VRIN matrix (\yes\ meets · \partly\ partly · \no\ does not)}\label{ex:vrin}
{\tablesetup
\begin{tabularx}{\linewidth}{P{38mm}LSSSSP{38mm}}
\toprule
\textbf{Resource / capability} & \textbf{Evidence that APSEZ possesses it} & \textbf{V} & \textbf{R} & \textbf{I} & \textbf{N} & \textbf{Implication; main vulnerability} \\
\midrule
R1 Mundra complex (concession, SEZ land, draft, JV terminals) & 45\% of domestic volume; 72\% margin; Mega Port; SEZ contiguous to port & \yes & \yes & \partly* & \partly & Sustained-advantage candidate; *tenure ends 16 Feb 2031 \\ \rowsep
R2 Non-Mundra domestic ports & 14 ports and terminals; margins 41--65\% & \yes & \partly & \no & \no & Parity; acquired in factor markets \\ \rowsep
R3 Transshipment (Vizhinjam, Colombo) & Vizhinjam ≈1.3m TEU FY26; Colombo Phase II 3.2m TEU & \yes & \yes & \partly & \partly & Temporary; state hubs, Colombo and Singapore rivals \\ \rowsep
R4 Shipping-line JV terminals & CT3 with MSC, CT4 with CMA CGM & \yes & \partly & \partly & \partly & Temporary; co-terminus with 2031 \\ \rowsep
R5 Rail, logistics parks, ICDs & 12 logistics parks; rail 695,517 TEU (FY26) & \yes & \no & \no & \no & Parity; CONCOR, public freight corridor \\ \rowsep
R6 Trucking and freight forwarding & Margins 8.4\% / 10.4\% (Q2 FY27) & \partly & \no & \no & \no & Parity or below \\ \rowsep
R7 Marine and dredging fleet & 249 vessels incl.\ 64 dredgers; marine margin 51\% & \yes & \partly & \partly & \partly & Temporary; captive value depends on R1 \\ \rowsep
R8 Project execution and acquisition integration & Capacity 653 → 671 MTPA; but capex overruns, Vizag exit & \yes & \partly & \partly & --- & Unproven \\ \rowsep
R9 Financial capacity & Net debt/EBITDA 1.8--1.9\texttimes; Baa3/BBB- ratings & \yes & \no & n/a & \no & Parity relative to sovereign-backed rivals \\ \rowsep
R10 Digital platforms & Vendor systems; ``proprietary'' trucking platform & \yes & \no & \no & \no & Parity (control case) \\
\bottomrule
\end{tabularx}}
\exsource{Sources: APSEZ (2026a; 2026b; 2026d); Shipping Tribune (2026).
Assessments are the author's.}
\end{exhibit}

%% =====================================================================
\section{VRIO analysis and comparison with VRIN}\label{sec:vrio}
%% =====================================================================
VRIO keeps value and rarity. It folds substitution into ``costly to imitate''
and adds a fourth question: is the firm organised to capture the value (Barney,
1995; Barney and Hesterly, 2021)? The evidence on organisation points both ways.

\emph{For} organisation, APSEZ has the following in place:
\begin{itemize}
  \item It reports return on capital for each segment, and from Q1 FY27 it
        reports results for each port cluster.
  \item It has kept to its leverage ceiling of 2.5 times EBITDA (actual
        1.8--1.9 times). All three international rating agencies moved their
        outlooks to stable or positive in FY26 (APSEZ, 2026a).
  \item It has built the ``shore-to-door'' bundle of port, rail, logistics
        parks and trucking.
  \item Its formal governance follows the listing framework. The chairman and
        the managing director come from the promoter family. Day-to-day
        management sits with a whole-time director and CEO. The audit committee
        is chaired by an independent director, and a separate reputation-risk
        committee exists (APSEZ, 2026c). The reputation-risk committee is an
        unusual structure that responds directly to the group's history of
        external scrutiny.
\end{itemize}

\emph{Against} organisation, the record shows three weaknesses:
\begin{itemize}
  \item \textbf{Weak capital planning.} FY26 capex overshot guidance, and FY27
        capex guidance was raised by ₹3,000 cr in mid-year (APSEZ, 2026d).
  \item \textbf{Contract-governance gaps.} The supplementary concession for
        Mundra's West Basin, operating since 2011, is still ``in the process of
        getting signed''. The sub-concessions for the MSC and CMA CGM terminals
        still await state approval (APSEZ, 2026b).
  \item \textbf{Related-party acquisitions.} NQXT was bought from a
        promoter-group company and paid for with new shares. Minority
        shareholders approved the deal (APSEZ, 2025). Over the same period,
        revenue from group companies rose by 72\%. In September 2025 the
        securities regulator, SEBI, found that the related-party allegations
        made by Hindenburg Research in 2023 were not established (SEBI, 2025).
        Even so, investors are likely to scrutinise this pattern.
\end{itemize}

Exhibit~\ref{ex:vrio} shows the resulting VRIO outcomes. The Mundra complex is a
sustained competitive advantage, but conditional. Operationally APSEZ is well
organised to exploit it. Contractually it is not, because unsigned agreements
and a fixed expiry date expose the asset that funds the group. Transshipment,
the joint ventures and the marine fleet give temporary advantages. Logistics,
rail, finance and digital systems give parity. Integration capability remains
unproven.

The two frameworks agree on the strongest resource but for different reasons,
and Exhibit~\ref{ex:diverge} shows where they diverge. VRIN's
non-substitutability test exposes an \emph{external} threat: a state-built
alternative bundle for containers. VRIO's organisation test exposes an
\emph{internal} one: contract governance and capital planning. The joint
ventures (R4) drop a level under VRIO because their legal basis is unsigned.
Integration capability (R8), which VRIN cannot resolve, looks weaker once
actual post-acquisition results are examined. Financial capacity fails VRIN
because sovereign capital substitutes for it, while VRIO rates it as parity.
VRIN therefore says more about the strength of the resource bundle, and VRIO
says more about whether management turns it into returns.

% ---------------------------------------------------------------- Exhibit 5
\begin{exhibit}[!htbp]
\caption{VRIO matrix and competitive outcome}\label{ex:vrio}
{\tablesetup
\begin{tabularx}{\linewidth}{P{30mm}SSP{18mm}LP{30mm}}
\toprule
\textbf{Resource} & \textbf{V} & \textbf{R} & \textbf{Costly to imitate} & \textbf{Organised (key evidence)} & \textbf{Outcome} \\
\midrule
R1 Mundra complex & \yes & \yes & ✓ until 2031 & ◐ Strong in operations; unsigned supplementary and sub-concessions & \textbf{Sustained, conditional on renewal} \\ \rowsep
R2 Non-Mundra ports & \yes & \partly & \no & ◐ Flat volume; east margin down & Parity \\ \rowsep
R3 Transshipment & \yes & \yes & \partly & ✓ Rapid ramp-up (Colombo 2m TEU in 18 months) & Temporary advantage \\ \rowsep
R4 Shipping-line JVs & \yes & \partly & \partly & ✗ Rights co-terminus with 2031, unsigned & Temporary advantage \\ \rowsep
R5 Rail / logistics parks & \yes & \no & --- & ✓ Bundled offer & Parity \\ \rowsep
R6 Trucking / forwarding & \partly & \no & --- & ◐ ROCE 10\% (segment) & Parity \\ \rowsep
R7 Marine fleet & \yes & \partly & \partly & ✓ Take-or-pay contracts & Temporary advantage \\ \rowsep
R8 Execution / integration & \yes & \partly & \partly & ◐ Overruns; Colombo slipped from Dec 2024 to Apr 2025 & Unproven \\ \rowsep
R9 Financial capacity & \yes & \no & --- & ✓ Leverage policy respected & Parity \\ \rowsep
R10 Digital & \yes & \no & --- & --- & Parity \\
\bottomrule
\end{tabularx}}
\exsource{Sources: APSEZ (2025; 2026a; 2026b; 2026d); Lloyd's List (2025);
Swarajya (2024). Outcome categories follow Barney and Hesterly (2021).}
\end{exhibit}

% ---------------------------------------------------------------- Exhibit 6
\begin{exhibit}[!htbp]
\caption{Where VRIN and VRIO diverge, and why}\label{ex:diverge}
{\tablesetup
\begin{tabularx}{\linewidth}{P{24mm}P{36mm}P{32mm}L}
\toprule
\textbf{Resource} & \textbf{VRIN verdict} & \textbf{VRIO verdict} & \textbf{Reason for the difference} \\
\midrule
R1 Mundra & Strong; N only partial & Sustained but conditional & N flags an external alternative bundle (JNPA + Vadhavan + rail corridor). O flags internal contract gaps. \\ \rowsep
R4 JVs & Partly rare, partly protected & Temporary & O fails: legal basis unsigned and expiring \\ \rowsep
R8 Integration & Indeterminate & Unproven & O tests actual post-deal outcomes \\ \rowsep
R9 Finance & Fails N & Parity & N names the substitute (sovereign capital); VRIO treats imitation and substitution together \\
\bottomrule
\end{tabularx}}
\end{exhibit}

%% =====================================================================
\section{Ansoff matrix: growth vectors and their evaluation}\label{sec:ansoff}
%% =====================================================================
Exhibit~\ref{ex:ansoff} maps APSEZ's announced initiatives, dated to its
disclosures up to 28 October 2026. It keeps them separate from this report's
proposals, which are shown in italics and developed in
Section~\ref{sec:recs}. Two classification choices need explaining. Entering a
new Indian hinterland, for example Gopalpur or Paradip, counts as market
development rather than penetration, because the customers and catchment are
new. Offshore and subsea vessel services count as related diversification. They
use marine skills, but they serve offshore-energy customers that APSEZ's ports
do not.

\textbf{Market penetration} at Mundra is the most suitable and acceptable
vector. It builds on the one sustained advantage and earns the highest returns.
Mundra's volume grew 8\% in H1 FY27. However, the timing is a problem. New
berths, such as the fifth container terminal and the VLCC jetty, will pay back
largely after February 2031, so their value depends on renewal terms that APSEZ
does not control.

\textbf{Market development} has two parts:
\begin{itemize}
  \item \textbf{In India}, it spreads concession risk across states. The returns
        so far are weak: non-Mundra volume is flat.
  \item \textbf{Abroad}, international ROCE was 8\% in FY26, excluding NQXT. The
        Q2 FY27 margin rise to 48\% reflects NQXT's contribution alongside the
        Colombo and Tanzania ramp-ups (APSEZ, 2026d). Because NQXT is a
        coal-export terminal, it also adds exposure to the energy transition.
        Haifa adds geopolitical exposure: APSEZ reported that the West Asia
        crisis disrupted logistics in Q1 FY27 (APSEZ, 2026d).
\end{itemize}

\textbf{Product development} (trucking, freight forwarding, warehousing, rail)
is suitable because it deepens the port bundle. Its acceptability varies.
Port-linked logistics parks and warehousing earn a 30\% EBITDA margin, but
trucking earns 8.4\%. \textbf{Related diversification} into offshore and subsea
services has the weakest link to APSEZ's sustained advantage and depends on
acquired capabilities. This is where the risk of over-diversification is
greatest.

The vectors also differ in the capital and execution they demand:
\begin{itemize}
  \item \textbf{Penetration} is the most capital-hungry. The Mundra expansion
        cleared in 2024 is a ₹45,000 cr programme (Business Standard, 2024), but
        it builds on proven operating routines.
  \item \textbf{International market development} uses large, lumpy deals.
        Colombo Phase II cost about US\$750m (APSEZ, 2026d). NQXT was paid for in
        new shares, which raised APSEZ's share count by about 6.7\% (APSEZ,
        2026a). Execution risk is political as well as operational.
  \item \textbf{Product development and diversification} need less capital per
        rupee of revenue but different management skills: managing trucking
        fleets, freight-forwarding sales and offshore-energy contracts. Most of
        these skills were acquired rather than built.
\end{itemize}

The danger is less that any single initiative fails than that management
attention and capital are spread across four vectors at once while the core
concession is still unresolved.

A rough cost-of-capital test sharpens these judgements. FY26 interest and bank
charges were about 7.6\% of average gross debt. That is the author's estimate
of the pre-tax cost of debt, not a WACC. Equity is costlier than debt, so
segments earning 8--10\% ROCE probably earn less than their full cost of
capital today. Management's FY31 targets of ₹91,500 cr revenue and ₹52,000 cr
EBITDA imply a 56.8\% margin, against 59\% in FY26 (APSEZ, 2026d). In effect,
management is planning for some dilution.

% ---------------------------------------------------------------- Exhibit 7
\begin{exhibit}[!htbp]
\caption{Ansoff matrix: existing initiatives (as of 28 Oct 2026) and
\prop{proposals in this report}}\label{ex:ansoff}
{\footnotesize
\renewcommand{\arraystretch}{1.2}%
\setlength{\tabcolsep}{5pt}%
\begin{tabularx}{\linewidth}{>{\centering\arraybackslash}m{7mm}|L|L}
 & \centering\textbf{Existing services} & \centering\arraybackslash\textbf{New services} \\
\hline
\rotatebox{90}{\textbf{Existing markets}} &
\textbf{Market penetration}
\begin{celllist}
  \item Mundra expansion: container terminal 5, VLCC jetty; ₹45,000 cr expansion cleared (2024)
  \item Empty-container yard at Mundra (Q2 FY27)
  \item Outcome: Mundra +8\% (H1 FY27), but national container share down 1.5 pts
  \item \prop{Proposed (Rec 3): multi-year contracts with container lines beyond the JV partners; rail-linked landed-cost offer}
\end{celllist} &
\textbf{Product development (port-linked)}
\begin{celllist}
  \item Rail (68 container rakes), 12 logistics parks, warehousing, agri-silos
  \item Trucking (3.4\texttimes\ revenue in FY26), international freight network
  \item Outcome: logistics ROCE 10\%; trucking margin 8.4\%
  \item \prop{Proposed (Rec 2): favour port-linked parks (30\% margin) over trucking}
\end{celllist} \\
\hline
\rotatebox{90}{\textbf{New markets}} &
\textbf{Market development}
\begin{celllist}
  \item India: Gopalpur, Paradip dry-bulk berths (30-yr, 18 MMT), Haldia, Vizhinjam export--import
  \item Abroad: Haifa, Dar es Salaam, Colombo Phase II (US\$750m), NQXT (acquired for shares, Dec 2025)
  \item Outcome: non-Mundra flat; international ROCE 8\% (excl.\ NQXT)
  \item \prop{Proposed (Rec 4): port-by-port review; sell minority stakes to recycle capital}
\end{celllist} &
\textbf{Related diversification}
\begin{celllist}
  \item Offshore support vessels (Astro Offshore, TAHID) in the Middle East, Africa and South Asia; subsea robotics; entry into the European subsea market (Q2 FY27)
  \item Standalone inland logistics parks (Kanpur, ₹1,500 cr)
  \item Outcome: marine ROCE 13\%
  \item \prop{Proposed (Rec 2): staged approval before further acquisitions}
\end{celllist} \\
\end{tabularx}}
\exsource{Sources: APSEZ (2025; 2026a; 2026d). Classification is the author's.}
\end{exhibit}

% ---------------------------------------------------------------- Exhibit 8
\begin{exhibit}[!htbp]
\caption{Evaluation of each growth vector}\label{ex:saf}
{\tablesetup
\begin{tabularx}{\linewidth}{P{25mm}LLLP{24mm}}
\toprule
\textbf{Vector} & \textbf{Suitability (fit with resources and forces)} & \textbf{Acceptability (return, risk)} & \textbf{Feasibility (capital, execution)} & \textbf{Verdict} \\
\midrule
Penetration (Mundra) & High: builds on R1; counters container rivalry & High return (23\% ROCE); renewal risk after 2031 & High: internal cash flow & Continue, gated on tenure \\ \rowsep
Market development, India & Medium: spreads grantor risk & Low so far: flat volume, east margin 41\% & High & Fix before expanding \\ \rowsep
Market development, international & Medium: hedge across geographies & Low--medium: 8\% ROCE; geopolitics; coal exposure & Medium: share-funded deals & Selective \\ \rowsep
Product development & High: deepens the bundle & Mixed: parks 30\% vs trucking 8\% margin & High: asset-light & Focus on port-linked services \\ \rowsep
Related diversification & Low: weak link to R1 & Medium: 13\% ROCE; acquired skills & Medium & Slow down; stage-gate \\
\bottomrule
\end{tabularx}}
\exsource{Criteria follow Whittington \textit{et al.} (2023). Data: APSEZ (2026a; 2026d).}
\end{exhibit}

%% =====================================================================
\section{Integrated strategic assessment}\label{sec:integration}
%% =====================================================================
Combining the frameworks produces one central finding. The resource that funds
APSEZ, the Mundra complex, is also the resource most exposed to the two
strongest external pressures. Those pressures are the grantor's bargaining power
as 2031 approaches and state-led container capacity on the west coast.
Meanwhile the Ansoff vectors mostly direct new capital to areas where VRIO finds
parity or temporary advantage. The strategy is therefore a hedge. It reduces
dependence on one site and one concession, but it dilutes returns. That
trade-off can be rational, but only if the new segments eventually earn above
their cost of capital, and the evidence does not yet show this.

The integration questions in the brief can be answered directly:
\begin{itemize}
  \item \textbf{Does APSEZ's VRIO position overcome the pressures identified by
        the Five Forces?} For bulk, liquids and cargo linked to the SEZ, largely
        yes. Site, land and dedicated infrastructure neutralise buyer power and
        deter entry. For containers, only partly: falling market share shows
        that the joint ventures and the logistics bundle have not offset rising
        rivalry.
  \item \textbf{Which resources best support growth?} Mundra supports
        penetration. The transshipment assets support market development. No
        sustained advantage clearly supports diversification.
  \item \textbf{Is diversification backed by organisational capability?} Marine
        diversification rests on acquired firms rather than on skills
        transferred from the ports, and trucking earns parity margins. So the
        organisational support is thinner than the growth targets imply.
  \item \textbf{Which external threats could erode the current advantages?}
        Three stand out: the energy transition (coal is about 30\% of domestic
        cargo, and NQXT is a coal terminal), conflict in West Asia (Haifa), and
        any shift in state policy on concession renewal.
\end{itemize}

The integration also resolves four apparent contradictions
(Exhibit~\ref{ex:integration}):
\begin{itemize}
  \item \textbf{Buyer power vs margins.} High container buyer power coexists
        with 72\% domestic margins because of cargo mix, captive harbour
        services and deconsolidated joint ventures. This is not evidence of weak
        buyers.
  \item \textbf{``Network'' vs portfolio.} The ``network'' that APSEZ describes
        behaves like a portfolio: when one site grows, the others do not.
  \item \textbf{Two kinds of discipline.} APSEZ is disciplined about leverage
        but not about capex planning.
  \item \textbf{Average tenure vs Mundra.} APSEZ reports an average remaining
        concession tenure of ``30+ years'' (APSEZ, 2026d). That average is not
        weighted by earnings, and Mundra, about a third of domestic-port EBITDA
        (31\% in Q2 FY27, before JV income), expires in just over four years.
\end{itemize}

Scale is best understood as a temporary advantage. Rivals can match capacity
within a decade. Vadhavan's planned first phase of about 15 million TEU
(Seatrade Maritime News, n.d.) is larger than APSEZ's current annualised
domestic container throughput of about 13 million TEU, although estimates of
when that phase will be complete range from around 2030 (ThePrint, 2024) to the
mid-2030s. What lasts is site quality, contiguous land and the time already
invested. Size alone does not.

Exhibit~\ref{ex:scenarios} sets out three scenarios for 2031:
\begin{itemize}
  \item \textbf{Renewal on similar terms.} The current strategy works and
        diversification is optional.
  \item \textbf{Renewal at a higher revenue share or with conditions.} Mundra's
        margin compresses, and the international and marine businesses become
        more valuable as earnings stabilisers.
  \item \textbf{Re-tendering or partial renewal.} This would be damaging. The
        container-line partners would gain bargaining power in the run-up,
        because their rights end on the same date.
\end{itemize}

An early Vadhavan opening would make the second and third scenarios worse.
APSEZ's position is therefore strong and likely to last until 2031. After 2031
it depends on two things: the concession outcome, and whether diversification
starts to earn its cost of capital.

% ---------------------------------------------------------------- Exhibit 9
\begin{exhibit}[!htbp]
\caption{Integration matrix: how the frameworks connect}\label{ex:integration}
{\tablesetup
\begin{tabularx}{\linewidth}{P{29mm}LP{22mm}LLP{21mm}}
\toprule
\textbf{Resource (VRIO outcome)} & \textbf{Forces it counters} & \textbf{Ansoff vectors it enables} & \textbf{Financial evidence} & \textbf{Threat that erodes it} & \textbf{Horizon} \\
\midrule
R1 Mundra (sustained, conditional) & Bulk/liquid buyers; entry & Penetration & Domestic ROCE 23\%; 72\% margin & Grantor power; Vadhavan & Secure to 2031; uncertain after \\ \rowsep
R3 Transshipment (temporary) & Rivalry for transshipment & Market development & International margin 48\% (Q2) & New state hubs; Colombo rivals & Medium \\ \rowsep
R4 JVs (temporary) & Container buyer power (partly) & Penetration & AICTPL margin $\sim$57\% & Partners' leverage near 2031 & To 2031 \\ \rowsep
R5--R6 Logistics (parity) & Little & Product development & ROCE 10\% & CONCOR; open corridor & --- \\ \rowsep
R7 Marine (temporary) & Supplier (dredging) power & Diversification & ROCE 13\% & Offshore cycle; acquired skills & Medium \\
\bottomrule
\end{tabularx}}
\end{exhibit}

% ---------------------------------------------------------------- Exhibit 10
\begin{exhibit}[H]
\caption{Strategy synthesis: pressures, resources, growth and responses}\label{ex:synthesis}
\centering
\begin{tikzpicture}[
    font=\fontsize{6.5}{8}\selectfont,
    execute at begin node={\hyphenpenalty=10000\exhyphenpenalty=10000},
    box/.style={draw=rulegrey, rounded corners=2.5pt, align=center,
                text width=30mm, minimum height=10.5mm, inner sep=2.5pt,
                line width=0.6pt},
    hd/.style={font=\bfseries\fontsize{7.4}{9}\selectfont,
               text=accent, align=center},
    link/.style={-{Stealth[length=4pt,width=3.2pt]}, draw=muted, line width=0.7pt}]
  % column headings
  \node[hd] at (17.5mm, 4mm)  {FIVE FORCES};
  \node[hd] at (59.0mm, 4mm)  {VRIO RESOURCES};
  \node[hd] at (100.5mm, 4mm) {ANSOFF VECTORS};
  \node[hd] at (141.0mm, 4mm) {RECOMMENDATIONS};
  % Five Forces pressures
  \node[box, fill=highbg] (f1) at (17.5mm, -7mm)  {Grantor power at 2031\\(supplier)};
  \node[box, fill=highbg] (f2) at (17.5mm, -21.5mm) {Container rivalry and\\state entry (Vadhavan)};
  \node[box, fill=modbg]  (f3) at (17.5mm, -36mm) {Container-line buyer power\\(co-opetition)};
  \node[box, fill=modbg]  (f4) at (17.5mm, -50.5mm) {East-coast rivalry;\\coal decline};
  % VRIO resources
  \node[box, fill=lowbg, draw=lowink, minimum height=15mm] (r1) at (59mm, -14mm)
       {\textbf{R1 Mundra complex}\\Sustained, conditional\\$\sim$⅓ of domestic EBITDA};
  \node[box, fill=modbg] (r2) at (59mm, -32.5mm) {R3/R4/R7 Temporary\\transshipment, JVs, marine};
  \node[box, fill=soft]  (r3) at (59mm, -48mm) {R2/R5/R6/R9 Parity\\other ports, logistics, finance};
  % Ansoff vectors
  \node[box, fill=lowbg]  (a1) at (100.5mm, -11.5mm) {Penetration\\ROCE 23\%};
  \node[box, fill=modbg]  (a2) at (100.5mm, -25.5mm) {Market development\\ROCE 8\% (intl)};
  \node[box, fill=modbg]  (a3) at (100.5mm, -39.5mm) {Product development\\ROCE 10\%};
  \node[box, fill=highbg] (a4) at (100.5mm, -53.5mm) {Diversification\\ROCE 13\%; weak link};
  % Recommendations
  \node[box, fill=white, draw=accent, text width=32mm] (c1) at (141mm, -7mm)
       {\textbf{1 Secure Mundra tenure}\\gate long-payback capex};
  \node[box, fill=white, draw=accent, text width=32mm] (c3) at (141mm, -21.5mm)
       {\textbf{3 Defend container share}\\before Vadhavan};
  \node[box, fill=white, draw=accent, text width=32mm] (c2) at (141mm, -36mm)
       {\textbf{2 Segment return gate}\\logistics, marine, intl};
  \node[box, fill=white, draw=accent, text width=32mm] (c4) at (141mm, -50.5mm)
       {\textbf{4 Turn around or recycle}\\non-Mundra ports};
  % links: forces -> resources
  \draw[link] (f1.east) -- (r1.west);
  \draw[link] (f2.east) -- (r1.west);
  \draw[link] (f3.east) -- (r2.west);
  \draw[link] (f4.east) -- (r3.west);
  % links: resources -> growth vectors
  \draw[link] (r1.east) -- (a1.west);
  \draw[link] (r2.east) -- (a2.west);
  \draw[link] (r2.east) -- (a4.west);
  \draw[link] (r3.east) -- (a3.west);
  % links: growth vectors -> recommendations
  \draw[link] (a1.east) -- (c1.west);
  \draw[link] (a1.east) -- (c3.west);
  \draw[link] (a2.east) -- (c4.west);
  \draw[link] (a3.east) -- (c2.west);
  \draw[link] (a4.east) -- (c2.west);
\end{tikzpicture}
\exsource{Colour shows strength: green = sustained / high return; amber =
temporary / moderate; red = high pressure or weak fit; grey = parity.}
\end{exhibit}

% ---------------------------------------------------------------- Exhibit 11
\begin{exhibit}[!htbp]
\caption{Mundra 2031 scenarios (qualitative, not forecasts)}\label{ex:scenarios}
{\tablesetup
\begin{tabularx}{\linewidth}{P{40mm}P{24mm}LP{38mm}}
\toprule
\textbf{Scenario} & \textbf{Effect on R1} & \textbf{Effect on strategy} & \textbf{Early warning indicator} \\
\midrule
A.\ Renewal on similar terms & Advantage extends & Current portfolio approach holds & Signed extension; signed sub-concessions \\ \rowsep
B.\ Renewal at higher revenue share or with conditions & Margin compresses & International and marine become stabilisers; return gate becomes critical & GMB terms published; royalty revision \\ \rowsep
C.\ Re-tender or partial renewal & Advantage at risk & JV partners' leverage peaks; Mundra capex stranded & Delay beyond 2028; tender notice \\
\bottomrule
\end{tabularx}}
\end{exhibit}

%% =====================================================================
\section{Strategic recommendations}\label{sec:recs}
%% =====================================================================
The four recommendations below follow from the preceding analysis, ranked by
impact and urgency (Exhibit~\ref{ex:priority}). Each draws on at least three
frameworks. Each is either new, or explicitly modifies an initiative APSEZ has
already announced. The KPI baselines come from the data used above.

\needspace{12\baselineskip}
\subsection{Recommendation 1 (priority 1): Make Mundra's tenure durable before 2029}
\begingroup\tablesetup
\begin{xltabular}{\linewidth}{P{30mm}L}
\toprule
\reclabel{Objective} & Turn the only sustained advantage from a contingent asset into a durable one, at least two years before the 16 February 2031 expiry. \\ \rowsep
\reclabel{Evidence} & Concession expires 2031; West Basin and the container JV sub-concessions are unsigned; Mundra carries 45\% of domestic volume (APSEZ, 2026b; 2026d). \\ \rowsep
\reclabel{Frameworks} & VRIO R1 conditional; supplier power (grantor); scenarios B and C; timing risk in Ansoff penetration. \\ \rowsep
\reclabel{Implementation} & (i) Finalise the outstanding supplementary and sub-concession agreements. (ii) Offer GMB a transparent value-sharing structure (for example, a revenue share that rises with volume) in exchange for tenure, accepting some margin. (iii) Negotiate renewal of the MSC and CMA CGM JV rights alongside multi-year volume commitments. (iv) Until tenure is secured, approve new Mundra projects only if they pay back before 2031. (v) Disclose an expiry schedule weighted by EBITDA. \textit{This modifies APSEZ's reported extension application (CargoConnect, n.d.).} \\ \rowsep
\reclabel{Resources and capital} & Mainly management and legal effort. The real cost is any higher royalty, so the design should limit the margin given up. \\ \rowsep
\reclabel{Risks} & Political sensitivity; terms worse than today; public scrutiny of any extension without competitive bidding. \\ \rowsep
\reclabel{Expected benefit} & Removes the largest single contingency on group value and protects about a third of domestic-port EBITDA. \\ \rowsep
\reclabel{KPIs} & Extension signed (target: by FY29); share of EBITDA with more than 10 years' tenure (baseline: Mundra, about 31\% of domestic EBITDA, has just over 4 years); share of Mundra capex with payback after 2031; JV volume commitments signed. \\
\bottomrule
\end{xltabular}
\endgroup

\needspace{12\baselineskip}
\subsection{Recommendation 2 (priority 2): Introduce a return gate for non-port segments}
\begingroup\tablesetup
\begin{xltabular}{\linewidth}{P{30mm}L}
\toprule
\reclabel{Objective} & Stop growth from diluting returns, by funding logistics, marine and international expansion only where it can beat the cost of capital. \\ \rowsep
\reclabel{Evidence} & Segment ROCE of 10\% (logistics), 13\% (marine) and 8\% (international) against 23\% for domestic ports; capex overshoot; implied FY31 margin of 56.8\% (APSEZ, 2026a; 2026d). \\ \rowsep
\reclabel{Frameworks} & VRIO parity for R5--R6; Ansoff evaluation; the 5F finding that logistics is competitive. \\ \rowsep
\reclabel{Implementation} & Publish a hurdle (for example, ROCE above the consolidated level within three years of investment). Stage-gate trucking and forwarding growth, and future marine acquisitions. Shift logistics capex towards port-linked parks and warehousing (30\% margin). Report ROCE for every segment including NQXT, and track capex against guidance. \\ \rowsep
\reclabel{Resources and capital} & Redirects part of FY27 capex (₹15--17k cr guidance); no new capital is needed. \\ \rowsep
\reclabel{Risks} & Slower revenue growth against the FY31 ambition; possible negative market reaction. \\ \rowsep
\reclabel{Expected benefit} & Holds consolidated ROCE near or above 16\% while the hedge is built. \\ \rowsep
\reclabel{KPIs} & Segment ROCE (FY26 baselines: 10\% / 13\% / 8\%); capex variance against guidance (baseline +28--39\%); share of capex going to segments below consolidated ROCE. \\
\bottomrule
\end{xltabular}
\endgroup

\needspace{12\baselineskip}
\subsection{Recommendation 3 (priority 3): Defend west-coast container share before Vadhavan}
\begingroup\tablesetup
\begin{xltabular}{\linewidth}{P{30mm}L}
\toprule
\reclabel{Objective} & Stop the decline in container share and lock in volume before new state capacity opens around 2030. \\ \rowsep
\reclabel{Evidence} & Container share fell from 45.9\% to 44.4\%; JNPA passed 8m TEU; APSEZ rail TEU fell 6.4\%; container buyer power is high (APSEZ, 2026d; Free Press Journal, 2026). \\ \rowsep
\reclabel{Frameworks} & 5F rivalry, entry and buyers; VRIN partial N for R1/R4; Ansoff penetration. \\ \rowsep
\reclabel{Implementation} & Multi-year service agreements with container lines beyond the two JV partners. A guaranteed door-to-door cost from Mundra to northern inland depots, combining port, rail and freight-corridor services. Use Vizhinjam as a feeder hub for APSEZ gateway ports. \textit{This focuses the existing shore-to-door strategy on containers.} \\ \rowsep
\reclabel{Resources and capital} & Moderate: rakes and terminal productivity. Not quantified from disclosed data. \\ \rowsep
\reclabel{Risks} & Price concessions erode margin; container lines may still split calls across ports. \\ \rowsep
\reclabel{Expected benefit} & Protects the contested segment, where buyer power is highest, before capacity arrives. \\ \rowsep
\reclabel{KPIs} & Container share (baseline 44.4\%, Q2 FY27); Mundra TEU growth against JNPA; rail TEU (baseline 167,453 per quarter); share of container volume under multi-year contracts (new disclosure). \\
\bottomrule
\end{xltabular}
\endgroup

\needspace{12\baselineskip}
\subsection{Recommendation 4 (priority 4): Turn around or recycle the non-Mundra portfolio}
\begingroup\tablesetup
\begin{xltabular}{\linewidth}{P{30mm}L}
\toprule
\reclabel{Objective} & Raise returns at ports with flat volumes and falling margins, or release the capital tied up in them. \\ \rowsep
\reclabel{Evidence} & Non-Mundra volume grew 0\% in H1 FY27; east-coast margin fell from 50\% to 41\%; Visakhapatnam coal terminal exit (APSEZ, 2026b; 2026d). \\ \rowsep
\reclabel{Frameworks} & VRIN R2 parity; factor-market logic (Barney, 1986); Ansoff domestic market development; east-coast rivalry. \\ \rowsep
\reclabel{Implementation} & A port-by-port review of cargo mix, with a planned move away from coal towards minerals and containers. Sell minority stakes in mature ports to long-term infrastructure investors, and redeploy the proceeds to R1 and R3. \textit{No such stake-sale programme was found in the disclosures reviewed.} \\ \rowsep
\reclabel{Resources and capital} & Releases rather than consumes capital; transaction costs. \\ \rowsep
\reclabel{Risks} & Valuation discounts; governance with co-investors; signalling. \\ \rowsep
\reclabel{Expected benefit} & Higher portfolio ROCE and funding for the sustained advantage, without extra leverage. \\ \rowsep
\reclabel{KPIs} & Non-Mundra volume growth (baseline 0\%); east-cluster EBITDA margin (baseline 41\%); capital recycled. \\
\bottomrule
\end{xltabular}
\endgroup

% ---------------------------------------------------------------- Exhibit 12
\begin{exhibit}[!htbp]
\caption{Prioritisation of recommendations}\label{ex:priority}
{\tablesetup
\begin{tabularx}{\linewidth}{P{5mm}LP{17mm}P{30mm}P{14mm}P{38mm}}
\toprule
\textbf{\#} & \textbf{Recommendation} & \textbf{Impact} & \textbf{Urgency} & \textbf{Ease} & \textbf{New or modified} \\
\midrule
1 & Durable Mundra tenure; capex gating & Very high & High (before 2029) & Low & Modifies existing application \\ \rowsep
2 & Segment return gate & High & High (FY27 capex raised) & High & New \\ \rowsep
3 & Container-share defence & High & Medium (before c.2030) & Medium & Focuses existing strategy \\ \rowsep
4 & Non-Mundra turnaround or recycling & Medium & Medium & Medium & New \\
\bottomrule
\end{tabularx}}
\end{exhibit}

%% =====================================================================
\section{Conclusion}\label{sec:conclusion}
%% =====================================================================
APSEZ is strongly placed in a moderately attractive industry, but its strength
is narrower and more time-bound than its scale suggests. The frameworks agree
that the Mundra complex is the real source of advantage. They disagree about its
durability, for reasons that are useful: the threats come from outside, through
state-led container capacity, and from inside, through contract governance.
Most of APSEZ's growth is placed in segments where it has parity, which is a
defensible hedge only if those segments start to earn their cost of capital.
The priority is therefore to secure Mundra before 2031 and to make the rest of
the portfolio earn its keep.

This assessment has limits. ROCE is company-defined. Peer data come from
secondary reports. No primary record of the concession negotiations was
available. Further research should obtain the Gujarat Maritime Board's terms,
the peers' audited accounts and an independent estimate of APSEZ's cost of
capital, since each could change the strength of the conclusions above.

%% =====================================================================
%%  References (Harvard, Cite Them Right)
%% =====================================================================
\clearpage
\phantomsection
\section*{References}
\addcontentsline{toc}{section}{References}
\label{sec:refs}
\begingroup
\small
\setlength{\parskip}{0pt}
\interlinepenalty=10000   % keep each reference entry on one page
\begin{list}{}{\setlength{\leftmargin}{8mm}\setlength{\itemindent}{-8mm}
               \setlength{\itemsep}{1.1ex}\setlength{\parsep}{0pt}
               \setlength{\topsep}{0pt}}
\raggedright

\item Adani Ports and Special Economic Zone Ltd (2025) \textit{APSEZ completes acquisition of NQXT Australia} [Press release]. 23 December. Available at: \url{https://www.adani.com/newsroom/media-releases/apsez-completes-acquisition-of-nqxt-australia} (Accessed: 4 November 2026).

\item Adani Ports and Special Economic Zone Ltd (2026a) \textit{Submission of media release and investor presentation on audited financial results for the quarter and financial year ended March 31, 2026} (Ref.\ APSEZL/SECT/2026-27/13). National Stock Exchange of India, 30 April. Available at: \url{https://nsearchives.nseindia.com/corporate/rkbhagia_30042026135819_IntimationofPRandIP.pdf} (Accessed: 4 November 2026).

\item Adani Ports and Special Economic Zone Ltd (2026b) \textit{Integrated annual report 2025-26: consolidated financial statements}. Ahmedabad: APSEZ. Available at: \url{https://www.adaniports.com/digital-reports/ports/2026/consolidated-financial-statements.pdf} (Accessed: 4 November 2026).

\item Adani Ports and Special Economic Zone Ltd (2026c) \textit{Integrated annual report 2025-26: statutory reports} (including Business Responsibility and Sustainability Report). Ahmedabad: APSEZ. Available at: \url{https://www.adaniports.com/digital-reports/ports/2026/statutory-reports.pdf} (Accessed: 4 November 2026).

\item Adani Ports and Special Economic Zone Ltd (2026d) \textit{Submission of media release and investor presentation on unaudited financial results for the quarter and half year ended September 30, 2026}. National Stock Exchange of India, 28 October. Available at: \url{https://nsearchives.nseindia.com/corporate/rkbhagia_28102026124435_IntimationofPRandIP.pdf} (Accessed: 4 November 2026).

\item Adler, N., Hirte, G., Kumar, S. and Niemeier, H.-M. (2022) `The impact of specialization, ownership, competition and regulation on efficiency: a case study of Indian seaports', \textit{Maritime Economics \& Logistics}, 24(3), pp.~507--536. doi: 10.1057/s41278-021-00200-y.

\item Ansoff, H.I. (1957) `Strategies for diversification', \textit{Harvard Business Review}, 35(5), pp.~113--124.

\item Barney, J.B. (1986) `Strategic factor markets: expectations, luck, and business strategy', \textit{Management Science}, 32(10), pp.~1231--1241. doi: 10.1287/mnsc.32.10.1231.

\item Barney, J. (1991) `Firm resources and sustained competitive advantage', \textit{Journal of Management}, 17(1), pp.~99--120.

\item Barney, J.B. (1995) `Looking inside for competitive advantage', \textit{Academy of Management Executive}, 9(4), pp.~49--61.

\item Barney, J.B. and Hesterly, W.S. (2021) \textit{Strategic management and competitive advantage: concepts and cases}. 6th edn. Hoboken, NJ: Pearson.

\item Brandenburger, A.M. and Nalebuff, B.J. (1996) \textit{Co-opetition}. New York: Currency Doubleday.

\item Business Standard (2024) `Adani gets environmental nod for Rs 45,000 cr Mundra port expansion', \textit{Business Standard}, 17 June. Available at: \url{https://www.business-standard.com/companies/news/adani-gets-environmental-nod-for-rs-45000-cr-mundra-port-expansion-124061700305_1.html} (Accessed: 4 November 2026).

\item CargoConnect (n.d.) `Gujarat at crossroads as pioneer port concessions near expiry without clear extension framework', \textit{Cargo Connect}. Available at: \url{https://cargoconnect.co.in/shipping/gujarat-at-crossroads-as-pioneer-port-concessions-near-expiry-without-clear-extension-framework} (Accessed: 4 November 2026).

\item Dierickx, I. and Cool, K. (1989) `Asset stock accumulation and sustainability of competitive advantage', \textit{Management Science}, 35(12), pp.~1504--1511. doi: 10.1287/mnsc.35.12.1504.

\item Free Press Journal (2026) `Navi Mumbai news: JNPA crosses 8 million TEUs, hits 100 MMT cargo in FY 2025--26, launches Jebel Ali--Jeddah service', \textit{Free Press Journal}. Available at: \url{https://www.freepressjournal.in/mumbai/navi-mumbai-news-jnpa-crosses-8-million-teus-hits-100-mmt-cargo-in-fy-202526-launches-jebel-alijeddah-service} (Accessed: 4 November 2026).

\item Government of India (2025) \textit{The Indian Ports Act, 2025} (Act No.~27 of 2025). India Code. Available at: \url{https://www.indiacode.nic.in/bitstream/123456789/22358/1/A2025-27.pdf} (Accessed: 4 November 2026).

\item Haralambides, H.E. (2002) `Competition, excess capacity, and the pricing of port infrastructure', \textit{International Journal of Maritime Economics}, 4(4), pp.~323--347. doi: 10.1057/palgrave.ijme.9100053.

\item India Shipping News (2021) `Concession agreements for 4 Gujarat ports to expire in 10 yrs; 3 operators seek extension', \textit{India Shipping News}. Available at: \url{https://indiashippingnews.com/?p=49390} (Accessed: 4 November 2026).

\item Lloyd's List (2025) `Adani's 3m teu Colombo terminal opens', \textit{Lloyd's List}, April. Available at: \url{https://www.lloydslist.com/LL1153129/Adani%E2%80%99s-3m-teu-Colombo-terminal-opens} (Accessed: 4 November 2026).

\item Maritime Gateway (2026) `Centre notifies India's first mega ports; Adani's Mundra among four to get new status', \textit{Maritime Gateway}, September. Available at: \url{https://www.maritimegateway.com/centre-notifies-indias-first-mega-ports-adanis-mundra-among-four-to-get-new-status/} (Accessed: 4 November 2026).

\item Notteboom, T.E. and Rodrigue, J.-P. (2005) `Port regionalization: towards a new phase in port development', \textit{Maritime Policy \& Management}, 32(3), pp.~297--313. doi: 10.1080/03088830500139885.

\item Paradip Port Authority (2026) \textit{Major ports record 915 million tonnes cargo in FY 2025-26, surpass target with 7.06\% growth}. 6 April. Available at: \url{https://paradipport.gov.in/all_news/major-ports-record-915-million-tonnes-cargo-in-fy-2025-26-surpass-target-with-7-06-growth/} (Accessed: 4 November 2026).

\item Porter, M.E. (2008) `The five competitive forces that shape strategy', \textit{Harvard Business Review}, 86(1), pp.~78--93.

\item Seatrade Maritime News (n.d.) `India to create \$9bn mega port near Mumbai', \textit{Seatrade Maritime News}. Available at: \url{https://www.seatrade-maritime.com/ports/india-create-9bn-mega-port-near-mumbai} (Accessed: 4 November 2026).

\item Securities and Exchange Board of India (SEBI) (2025) \textit{Final order in the matter of Adicorp Enterprises Private Limited}. September. Available at: \url{https://www.sebi.gov.in/sebi_data/attachdocs/sep-2025/order_matter_adicorp.pdf} (Accessed: 4 November 2026).

\item Shipping Tribune (2026) `Adani's Vizhinjam transhipment port handles 1.3 million TEUs in FY26', \textit{Shipping Tribune}. Available at: \url{https://www.shippingtribune.com/index.php/news/shipping/Adani%27s+Vizhinjam+transhipment+port+handles+1.3+million+TEUs+in+FY26} (Accessed: 4 November 2026).

\item Swarajya (2024) `Adani's Colombo port project: first phase expected to be operational by Dec 2024', \textit{Swarajya}. Available at: \url{https://swarajyamag.com/amp/story/infrastructure/adanis-colombo-port-project-first-phase-expected-to-be-operational-by-dec-2024} (Accessed: 4 November 2026).

\item ThePrint (2024) `Boost for maritime sector as Rs 76,220 cr Vadhavan port planned in Maharashtra gets cabinet nod', \textit{ThePrint}, June. Available at: \url{https://theprint.in/india/boost-for-maritime-sector-as-rs-76220-cr-vadhavan-port-planned-in-maharashtra-gets-cabinet-nod/2138626} (Accessed: 4 November 2026).

\item Whalesbook (2026) `JSW Infrastructure posts ₹5,361 Cr FY26 revenue, targets 400 mtpa capacity', \textit{Whalesbook}, May. Available at: \url{https://www.whalesbook.com/corporate-news/English/industrial-goodsservices/JSW-Infrastructure-Posts-indian-rupee5361-Cr-FY26-Revenue-Targets-400-mtpa-Capacity/6a01f1b9707d23e8445cdedc} (Accessed: 4 November 2026).

\item Whittington, R., Regn\'er, P., Angwin, D., Johnson, G. and Scholes, K. (2023) \textit{Exploring strategy: text and cases}. 13th edn. Harlow: Pearson.
\end{list}
\endgroup

%% =====================================================================
%%  Appendices
%% =====================================================================
\clearpage
\phantomsection
\section*{Appendix A: Verification and compliance checklist}
\addcontentsline{toc}{section}{Appendix A: Verification and compliance checklist}
\label{app:checklist}
\begingroup\tablesetup
\begin{xltabular}{\linewidth}{P{25mm}LLP{30mm}}
\toprule
\textbf{Review cycle} & \textbf{What was checked} & \textbf{Corrections made} & \textbf{Status} \\
\midrule
\endhead
1 Requirement coverage & All four frameworks; integration; 3--5 recommendations with 8 components; 7 visual types; executive summary to references & Ansoff evaluation criteria added; 2031 scenarios added & Satisfied \\ \rowsep
2 Evidence verification & All APSEZ figures checked against NSE filings and audited statements & Revenue basis difference (audited FY25 ₹30,475 cr vs restated ₹31,079 cr) footnoted; Q2 cargo corrected from a secondary 142 MMT (total) to the primary 118 MMT domestic / 24 MMT international split; claims narrowed to what sources support (NQXT valuation dropped; Vadhavan timing given as a range; international-margin drivers attributed as APSEZ states them); every in-text citation matched to the reference list & Satisfied for company data; peer data secondary (\hyperref[app:data]{Appendix~B}) \\ \rowsep
3 Framework accuracy & VRIN-N kept distinct from Porter's substitutes; substitution placed inside VRIO imitability; government treated as acting through the forces, with the grantor adaptation justified & Domestic new-catchment entry reclassified to market development; offshore/subsea reclassified to related diversification & Satisfied \\ \rowsep
4 Critical analysis & Each strength tested against counter-evidence (non-Mundra flat; JV deconsolidation; ``30+ years'' average) & ``Network'' claim downgraded to portfolio; acquisition capability rated unproven & Satisfied \\ \rowsep
5 Strategic consistency & Every recommendation traced to at least three frameworks and to KPI baselines & Early concession recommendation re-derived from evidence rather than assumed & Satisfied \\ \rowsep
6 Final audit & Word count; Harvard format; units and periods & See word count below & Main body 4,977 words (target 5,000) \\
\bottomrule
\end{xltabular}
\endgroup

{\footnotesize
Word count. Main body (Sections 1--11: prose, headings and recommendation
tables): 4,977 words, of which recommendation tables 730. Exhibits (tables,
diagrams, captions and source notes) add 1,647 words and are excluded, as are
references and appendices.\par}

\vspace{4mm}
\phantomsection
\section*{Appendix B: Data notes and limitations}
\addcontentsline{toc}{section}{Appendix B: Data notes and limitations}
\label{app:data}
{\footnotesize
\begin{itemize}
  \item \textbf{Basis.} Segment, EBITDA and ROCE figures use APSEZ's management
        basis (APSEZ, 2026a; 2026d). ROCE is company-defined and was not
        recomputed. The FY26 consolidated figure excludes NQXT.
  \item \textbf{Scope of derived ratios.} Utilisation (451 MMT \textdiv{} 653
        MTPA) uses year-end capacity, so it understates utilisation of average
        capacity. Revenue per tonne is not quoted as a tariff, because the
        domestic segment includes harbour and other revenue. The cost-of-debt
        proxy (≈7.6\%) includes bank charges and approximates FY25 gross debt
        from borrowings. JV figures (AICTPL) are on a 100\% basis and are not
        consolidated.
  \item \textbf{Market share.} Shares are calculated by APSEZ. The author's
        check (451 \textdiv{} 1,668 = 27.0\%) uses an all-India total that APSEZ
        compiled. The full-year FY26 government figure for non-major ports was
        not available.
  \item \textbf{Peers.} JNPA, JSW Infrastructure and Vadhavan figures come from
        press reports of those organisations' data. Their primary filings were
        not reviewed.
  \item \textbf{Not verified.} No primary GMB document on the Mundra extension
        was found. A report of a ``letter of comfort'' (August 2026) is
        single-source and is not relied on. The terms of the 2026 US case
        outcomes were not reviewed, and are not relied on.
  \item \textbf{Access.} Some trade-press sites were blocked by the research
        environment, and paywalled academic databases were not accessed.
        Academic references were checked against bibliographic records
        (RePEc/DOI), except the page range for Porter (2008).
  \item \textbf{Traceability.} The source register, data log and analysis notes
        are kept in the project repository (folder \texttt{research/}).
\end{itemize}}

\end{document}
